\chapter{Diagrama de casos de uso}
\label{cap:refactorizacion-diagrama}

Para que los nombres y las relaciones se lean con claridad, los casos de uso
se presentan en dos diagramas PlantUML independientes: gestión y
configuración en la figura~\ref{fig:casos-uso-gestion}; participación y
verificación en la figura~\ref{fig:casos-uso-votacion}. Las fuentes editables
se guardan en \texttt{Figuras/diagramas/}.

\begin{figure}[p]
  \centering
  \resizebox{0.94\textwidth}{!}{% Diagrama vectorial simplificado para lectura en la memoria.
\begin{tikzpicture}[font=\sffamily\scriptsize,
  actor/.style={draw=blue!55!black,fill=blue!4,rounded corners=2pt,align=center,text width=2.0cm,minimum height=.62cm,inner sep=3pt},
  uc/.style={draw=gray!65!black,fill=blue!3,ellipse,align=center,text width=3.0cm,minimum height=.72cm,inner sep=2pt},
  rel/.style={draw=gray!70,thin},
  gen/.style={draw=gray!70,thin,-{Triangle[open]}}]
\node[font=\bfseries\small] at (4,5.8) {Casos de uso: gestión y configuración};
\draw[rounded corners,draw=gray!55] (0.5,-5.5) rectangle (7.6,5.35);
\node[fill=white,font=\bfseries] at (4,5.35) {Plataforma SaaS (Gestión)};

\node[uc] (rf01) at (2.5,4.55) {RF-01\\Organizaciones};
\node[uc] (rf02) at (2.5,3.55) {RF-02\\Roles y permisos};
\node[uc] (rf03) at (2.5,2.55) {RF-03\\Gestionar censo};
\node[uc] (rf04) at (2.5,1.55) {RF-04\\Autenticar gestión};
\node[uc] (rf05) at (2.5,.55) {RF-05a\\Iniciar sesión gestión};
\node[uc] (rf06) at (2.5,-.45) {RF-06\\Configurar proceso};
\node[uc] (rf07) at (2.5,-1.45) {RF-07\\Ciclo de vida};
\node[uc] (rf14) at (2.5,-2.45) {RF-14\\Consultar participación};
\node[uc] (rf16) at (2.5,-3.45) {RF-16\\Generar acta PDF};
\node[uc] (rf17) at (2.5,-4.45) {RF-17\\Enviar notificaciones};
\node[uc] (rf18) at (5.8,3.55) {RF-18\\Privacidad RGPD};
\node[uc] (rf19) at (5.8,1.55) {RF-19\\Trazabilidad};
\node[uc] (rf13a) at (5.8,-.45) {RF-13a\\Escrutinio al cierre};

\node[actor] (super) at (-2,4.55) {Superadministrador};
\node[actor] (gestion) at (-2,1.0) {Usuario de gestión};
\node[actor] (censo) at (-2,-1.45) {Gestor de censo};
\node[actor] (aud) at (-2,-3.5) {Usuario con auditoría};
\node[actor] (admin) at (-5,-4.55) {Administrador};
\node[actor] (auditor) at (-5,-2.45) {Auditor};

\begin{pgfonlayer}{background}
\draw[rel] (super.east)--(rf01.west);
\foreach \u in {rf02,rf04,rf05,rf06,rf07,rf14,rf16,rf17,rf18} \draw[rel] (gestion.east)--(\u.west);
\draw[rel] (censo.east)--(rf03.west);
\draw[rel] (aud.east)--(rf19.west);
\draw[dashed,gray] (rf07.east)--node[above,font=\tiny]{\textless\textless include\textgreater\textgreater}(rf13a.west);
\draw[gen] (super.south)--(gestion.north);
\draw[gen] (censo.north)--(gestion.south);
\draw[gen] (aud.north)--(gestion.south);
\draw[gen] (admin.north)--(aud.south);
\draw[gen] (auditor.south)--(aud.north);
\end{pgfonlayer}
\node[font=\tiny,align=left] at (4,-5.05) {Las especializaciones heredan las asociaciones de sus actores genéricos.};
\end{tikzpicture}
}
  \caption{Diagrama de casos de uso (configuración y administración).}
  \label{fig:casos-uso-gestion}
\end{figure}

\begin{figure}[p]
  \centering
  \resizebox{0.94\textwidth}{!}{% Diagrama vectorial simplificado para lectura en la memoria.
\begin{tikzpicture}[font=\sffamily\scriptsize,
  actor/.style={draw=blue!55!black,fill=blue!4,rounded corners=2pt,align=center,text width=2.0cm,minimum height=.62cm,inner sep=3pt},
  uc/.style={draw=gray!65!black,fill=blue!3,ellipse,align=center,text width=3.5cm,minimum height=.74cm,inner sep=2pt},
  rel/.style={draw=gray!70,thin}]
\node[font=\bfseries\small] at (3.5,8.65) {Casos de uso: participación y verificación};
\draw[rounded corners,draw=gray!55] (0,.2) rectangle (7.1,7.35);
\node[fill=white,font=\bfseries] at (3.55,7.35) {Plataforma SaaS (Votación)};

\node[uc] (rf05) at (3.55,6.8) {RF-05b\\Iniciar sesión de votante};
\node[uc] (rf08) at (3.55,5.85) {RF-08\\Acceder como votante};
\node[uc] (rf09) at (3.55,4.9) {RF-09\\Emitir voto};
\node[uc] (rf10) at (3.55,3.95) {RF-10\\Obtener y verificar recibo};
\node[uc] (rf12) at (3.55,3.0) {RF-12\\Auditoría hash-chain};
\node[uc] (rf13b) at (3.55,2.05) {RF-13b\\Verificar resultado};
\node[uc] (rf15) at (3.55,.92) {RF-15\\Consultar/exportar resultados};

\node[actor] (votante) at (-2,5.7) {Votante};
\node[actor] (auditor) at (-2,3.0) {Auditor};
\node[actor] (admin) at (-2,.92) {Administrador};

\begin{pgfonlayer}{background}
\draw[rel] (votante.east)--(rf05.west);
\draw[rel] (votante.east)--(rf08.west);
\draw[rel] (votante.east)--(rf09.west);
\draw[rel] (votante.east)--(rf13b.west);
\draw[rel] (auditor.east)--(rf12.west);
\draw[rel] (admin.east)--(rf15.west);
\draw[dashed,-{Triangle[open]}] (rf10.east) -- ++(.65,0) |- node[pos=.25,right,font=\tiny]{\textless\textless extend\textgreater\textgreater} (rf09.east);
\end{pgfonlayer}
\node[draw=orange!65!black,fill=orange!7,rounded corners=2pt,align=center,text width=6.2cm,inner sep=4pt] at (3.55,-.55) {Punto de extensión RF-09: tras confirmar la emisión. RF-10 es opcional y no tiene asociación directa con Votante.};
\end{tikzpicture}
}
  \caption{Diagrama de casos de uso (participación y verificación).}
  \label{fig:casos-uso-votacion}
\end{figure}

El modelo usa asociaciones sin punta de flecha entre actores y casos de uso.
La generalización de actores representa especialización de roles; por ello,
Superadministrador y Gestor de censo heredan de \emph{Usuario de gestión},
mientras Administrador y Auditor heredan además de \emph{Usuario con auditoría}.
Este actor intermedio mantiene RF-19 limitado a esos dos perfiles. Cada caso de uso conserva como máximo un actor asociado directamente.

El requisito RF-05, «Iniciar sesión», se divide en dos casos trazables al requisito de origen: RF-05a para los perfiles de gestión y RF-05b para el votante. OIDC/SSO se conserva como decisión de diseño y mecanismo posible; no forma parte del nombre genérico del requisito. La integración del proveedor de identidad se documenta como decisión arquitectónica y no añade una segunda asociación directa a los casos de uso. RF-05b permanece separado del acceso a la papeleta (RF-08). El caso RF-13 se desglosa en RF-13a, incluido desde RF-07 para el
escrutinio obligatorio al cierre, y RF-13b, que permite verificar el resultado
publicado.

RF-10 y RF-11 se modelan como extensiones opcionales de RF-09, sin asociación
directa con el votante. El punto de extensión de RF-09 se sitúa tras confirmar
la emisión. RF-19 se asocia a \emph{Usuario con auditoría}; los perfiles
Administrador y Auditor lo heredan.
